\begin{figure}[htbp]
    \centering
    \begin{tikzpicture}[
        node distance=2cm,
        block/.style={
            rectangle,
            minimum width=2cm,
            % minimum height=1.5cm,
            % text width=4cm,
            align=center,
            rounded corners=4pt,
            draw,
            font=\small
        },
        ancestry/.style={
            fill=Mahogany!10,
        },
        newick/.style={
            fill=Peach!10,
        },
        arrow/.style={
            -Latex,
            thick,
            bend left=45
        },
        title/.style={
            font=\bfseries\small, % Bold and small font for titles
            anchor=south, % Anchor the title to the bottom of the node to align with the top row
        },
        scale=0.9, every node/.style={transform shape}
    ]

    % Define nodes for each step
    \node (step0) [newick, block] {`(\textcolor{red}{(0,2)6},1)7,(3,4)5)8;'};

    \node (step1anc) [ancestry, block, below right=0.25cm and 1.5cm of step0] {$\begin{pmatrix} 0 & 2 & 6 \end{pmatrix}$};
    \node (step1newick) [newick, block, below left=0.25cm and 1.5cm of step1anc] {`((\textcolor{red}{(}\textcolor{Purple}{6}\textcolor{red}{,1)7},(3,4)5)8;'};
    
    \node (step2anc) [ancestry, block, below right=0.25cm and 1.5cm of step1newick] {$\begin{pmatrix} 0 & 2 & 6 \\ 6 & 1 & 7 \end{pmatrix}$};
    \node (step2newick) [newick, block, below left=0.25cm and 1.5cm of step2anc] {`(\textcolor{Purple}{7},\textcolor{red}{(3,4)5})8;'};

    \node (step3anc) [ancestry, block, below right=0.25cm and 1.5cm of step2newick] {$\begin{pmatrix} 0 & 2 & 6 \\ 6 & 1 & 7 \\ 3 & 4 & 5 \end{pmatrix}$};
    \node (step3newick) [newick, block, below left=0.25cm and 1.5cm of step3anc] {`\textcolor{red}{(7,\textcolor{Purple}{5})8};'};
    
    \node (step4anc) [ancestry, block, below right=0.25cm and 1.5cm of step3newick] {$\begin{pmatrix} 0 & 2 & 6 \\ 6 & 1 & 7 \\ 3 & 4 & 5 \\ 7 & 5 & 8 \end{pmatrix}$};
    \node (step4newick) [newick, block, below left=0.25cm and 1.5cm of step4anc] {`\textcolor{Purple}{8};'};

    % Title nodes
    \node (newicktitle) [title, above=0.25cm of step0] {Newick string};
    \node (ancestrytitle) [title, above=0.25cm of step1anc] {Ancestry};

    % Draw the "snake" arrows
    \draw [arrow] (step0.south east) to (step1anc.north west);
    \draw [arrow, Purple!70] (step1anc.south west) to (step1newick.north east);
    \draw [arrow] (step1newick.south east) to (step2anc.north west);
    \draw [arrow, Purple!70] (step2anc.south west) to (step2newick.north east);
    \draw [arrow] (step2newick.south east) to (step3anc.north west);
    \draw [arrow, Purple!70] (step3anc.south west) to (step3newick.north east);
    \draw [arrow] (step3newick.south east) to (step4anc.north west);
    \draw [arrow, Purple!70] (step4anc.south west) to (step4newick.north east);

    % Draw annotation arrows
    \node (extra1) [above=1cm of ancestrytitle, xshift=3cm] {Add triplet to ancestry};
    \draw[-Latex, thick, dashed] 
        (extra1.south west) 
        -- ($($(step0.south east)!0.33!(step1anc.north west)$) + (0.25cm,0.35cm)$);

    \node (extra2) [above=0.1cm of step2anc, xshift=3cm, align=center, text=Purple] {Keep the parental node \\ and find the next triplet};
    \draw[-Latex, thick, dashed, Purple] 
        (extra2.west) 
        -- ($($(step1anc.south west)!0.4!(step1newick.north east)$) + (0.25cm,-0.25cm)$);

    \end{tikzpicture}
    \caption{Illustration of the parsing from a Newick string to the ``ancestry" object (also referred as the ``ancestry" matrix in Scheidwasser \textit{et al.}~\cite{scheidwasser2025}). In the absence of parental nodes, the third element of the triplet is assigned as the maximum of the two first elements, and we keep the minimum.}
    \label{fig:parse-snake}
\end{figure}